\subsection{Clifford 部分から外れる中心三次 incidence の一次方向}
\label{subsec:clifford-transverse-tangent}
前節の全切断分類は、最高次の反交換関係を固定した部分の分類である。
その外への候補を、中心三次元が持ち上がるという有限次数の条件から調べる。
以下は任意のよい $h$ と全ての $\beta,A,B,C,D$ に対する一次の計算であり、
形式変形の積分可能性や新しい H3 成分の存在を結論しない。

$k_\epsilon=k[\epsilon]/(\epsilon^2)$ とし、中心 $t$ を保つ関係式を
\begin{equation}\label{eq:clifford-tangent-relations}
 (1+\epsilon\gamma)yz+(1-\epsilon\gamma)zy
 =hx^2+\epsilon\alpha_x tx+\beta_x t^2
\end{equation}
およびその巡回置換とする。$h,\beta$ はこの式では固定する。
切断の特殊繊維は前節と同じ
\[
 G_0=g_\beta+t(Ax^2+By^2+Cz^2)+Dt^3
\]
であり、$\dot G$ には任意の同次三次元を許す。

\begin{proposition}[指定四方向の中心 incidence 接空間]
\label{prop:clifford-transverse-tangent}
関係式 \eqref{eq:clifford-tangent-relations} の下で
$G_0+\epsilon\dot G$ が中心となるための必要十分条件は
\begin{equation}\label{eq:clifford-tangent-criterion}
 \alpha_x=-4A\gamma,\qquad
 \alpha_y=-4B\gamma,\qquad
 \alpha_z=-4C\gamma
\end{equation}
である。固定したこの包絡方向について、持ち上げの全体は
$Z(B_{h,\beta})_3$ を平行移動した五次元アフィン空間となる。
特に $\gamma=1$ に対する一つの持ち上げは、$d=1+h^3$ として
\begin{align*}
 d\dot G={}&-3hy^3+(h^3-2)zyx
       +2h^2t(Cyx-Ayz-Bzx)\\
 &-t^2\bigl((4h^2A^2+4hBC-\beta_x)x
        +2(2hAC+2h^2B^2+\beta_y)y\\
 &\hspace{37mm}+(4hAB+4h^2C^2-\beta_z)z\bigr).
\end{align*}
全ての係数で次数 $3,4$ の包絡成分も $k_\epsilon$ 上自由であり、
その階数はそれぞれ $20,35$ である。
\end{proposition}

\begin{proof}
必要な線形代数を、特殊点でも有効な単元小行列式とともに記す。
中心 $t$ を左に移し、次数 $n$ の単語 $t^{n-|w|}w$ を $|w|\leq n$ の
活動変数の単語で表す。これは $t=1$ という記法による有限次元計算であり、
$t$ を可逆化して低次数の関係を推測する操作ではない。
係数環を
\[
 \mathcal A=\mathbb Q[h,\beta_x,\beta_y,\beta_z,A,B,C,D,
                       h^{-1},(1+h^3)^{-1}]
\]
とする。特殊繊維の三関係式に左右から長さの和が $2$ 以下の単語を
挿入した行列を $M$ とすると、その大きさは $121\times102$ である。
$86$ 本の独立列からなる $N$、関係係数行列 $T$、商への射影 $P$ を
\[
 M=NT,\qquad PN=0,\qquad \operatorname{rank}P=35
\]
となるように取れる。$N$ の証明書の小行列式は
$h^{14}(1+h^3)^{11}$ であり、従って全ての許された係数で単元である。
四つの包絡方向の関係行列の微分を $M_j$、選んだ列の微分を $N_j$ とすると
\[
 P(M_j-N_jT)=0\qquad (j=\gamma,\alpha_x,\alpha_y,\alpha_z)
\]
が係数環全体で成り立つ。これは選んだ関係の間の全ての syzygy も
一次に持ち上がることを表し、次数 $4$ の自由性を与える。
次数 $3$ にも同様に単元行列式を持つ $20$ 語の商基底を取る。

この基底の一般の補正 $\dot G$ の三つの交換子を商に射影した
$105\times20$ 行列を $Z$ とする。$Z$ は階数 $15$ で、選んだ
$15\times15$ 小行列式は $-256h^6$ となる。
その像への射影から、$EZ=0$ となる多項式係数の行列 $E$ を取れる。
$[G_0,x_i]=Nc_i$ と書き、四方向について $-PN_jc_i$ を並べた
$105\times4$ 行列を $V$ とすると、中心性の方程式は
\[
 Z\dot G+V(\gamma,\alpha_x,\alpha_y,\alpha_z)^T=0.
\]
$EV$ の三つの行は正確に
\[
 \begin{pmatrix}
 -4A/h&-1/h&0&0\\
 6B/d&0&3/(2d)&0\\
 6Ch/d&0&0&3h/(2d)
 \end{pmatrix}
\]
となる。後三列の行列式は $-9/(4d^2)$ で単元だから、
\eqref{eq:clifford-tangent-criterion} が必要である。
逆に上に表示した $\dot G$ を代入すると、$\gamma=1$ と
$\alpha=-4(A,B,C)$ に対して全 $105$ 方程式が成り立つ。
一般の $\gamma$ はこの持ち上げを $\gamma$ 倍すればよい。
解の差は $\ker Z=Z(B_{h,\beta})_3$ であり、前節より五次元である。

以上の全行列恒等式と単元行列式は
\texttt{verify\_cubic\_clifford\_tangent.py} が $\mathcal A$ 上で検算する。
有理関数体で計算しただけで終えず、使用する逆行列・射影・持ち上げの
分母が $h$ と $1+h^3$ の因子に限られることも確認する。
従って有限個の標本からの外挿も、$\beta_i\ne0$ という追加条件もない。
\end{proof}

$\gamma\ne0$ は反交換関係の最高次部分を変えるので、前節の指定
Clifford 表示の部分から外れる一次 incidence 方向が得られた。
ただしこの命題は、全ての包絡変形方向の計算、二次以上の障害の消滅、
積分した切断の H3 条件、あるいは七次元商が全 moduli の成分でないという
断定を与えない。以下ではこの明示式による第二次の延長を示す。
その先の中心性の形式積分は次節で解決する。実際の H3 族への実現は別の問題である。

\begin{corollary}[指定方向の二次障害の消滅]
\label{cor:clifford-transverse-second}
上の明示的 $\dot G$ は、全係数に対して $\epsilon^3=0$ まで延長できる。
具体的には $\gamma$ の項を $\epsilon$ のまま保ち、対角係数を
\[
 -4\epsilon(A,B,C)+\epsilon^2\kappa,
 \qquad
 \kappa=\frac4{1+h^3}
 \bigl(A(h^3+3),\,-B(3h^3+1),\,-C(h^3-3)\bigr)
\]
とする。このとき中心切断
$G_0+\epsilon\dot G+\epsilon^2G_2$ が存在し、次数 $4$ 成分はなお自由である。
\end{corollary}
\begin{proof}
一次方向を代入した関係行列を $M_1$、選んだ独立列を $N_1$ とする。
前の証明の単元小行列から
\[
 NT_1=M_1-N_1T,\qquad PN_1T_1=0
\]
を得る。第二式は、第二次の新しい包絡係数をまだ加えない場合の
全 syzygy 障害が消えることを表す。
$[G_0,x_i]=Nc_{0i}$ とし、一次の式を
$[\dot G,x_i]-N_1c_{0i}=Nc_{1i}$ と解く。
第二次の中心性の残差は $H_i=-PN_1c_{1i}$ である。
三交換子についてこれを並べた $H$ に対し、前の三つの単元係数の
障害行を使って第二次の対角係数を解くと、表示した $\kappa$ を得る。
残りの全行についても
\[
 E(H+V_\alpha\kappa)=0
\]
が成り立つ。従って $ZG_2+H+V_\alpha\kappa=0$ を解ける。
$M_2,N_2$ をこの対角係数に対応する行列とすると、さらに
\[
 P(M_2-N_1T_1-N_2T)=0
\]
であり、全関係の第二次 syzygy も持ち上がる。
\texttt{verify\_cubic\_clifford\_second\_order.py} はこれらの全行列式を
同じ係数環上で検証し、$G_2$ の全八項と分母の監査を証明書へ保存する。
これは $\epsilon^3=0$ の主張であり、全次数の形式積分を意味しない。
\end{proof}
